\section{Ad Hoc Teamwork：和陌生队友临场组队}

从人类反馈转向多体协作后，训练分布又多了一层未知性：部署时的伙伴可能从未出现过。本节回答 Agent 如何在不重训队友、不共享私有协议的情况下，凭少量在线行为推断角色并调整自己的贡献。

Stone 把 ad hoc teamwork 作为 interaction 方向的代表问题：目标是让 agent 在\textbf{无法预先控制队友策略}的情况下，仍能快速形成高质量协作。

\begin{figure}[H]
\centering
\includegraphics[width=0.95\textwidth]{diagrams/ad-hoc-teamwork.png}
\caption{讲义整理图：观察陌生队友、推断角色并在线适配}
\label{fig:ad-hoc-teamwork}
\end{figure}

\begin{knowledgebox}{概念解释：Ad Hoc Teamwork 的最小闭环}
Agent 先观察伙伴动作和通信，再更新对角色、能力与意图的 belief，随后调整自己的策略；共同结果又成为新证据。关键约束是不能要求伙伴重新训练或遵守私有协议。LLM 多 Agent 若只在固定 prompt 和同模型副本间测试，尚未证明具备这种能力。
\end{knowledgebox}

\begin{importantbox}{问题定义}
Ad hoc teamwork 的关键约束是：
\begin{itemize}
    \item 队友模型不可预训练到固定分布；
    \item 协同策略必须在线适配；
    \item 任务成功依赖 ``对队友行为的实时解释''。
\end{itemize}
\end{importantbox}

这和今天多 Agent coding、multi-role planner、企业流程自动化高度一致。你经常要和 ``未知策略的外部参与者'' 共事，比如新接入 tool、异构模型、甚至人类操作员。

\begin{knowledgebox}{对 LLM Multi-Agent 的直接启发}
如果系统只在固定队友配置下验证，部署后经常会遇到协同退化。更稳健的方法是把 ``队友分布扰动'' 作为训练和评测的一部分。
\end{knowledgebox}

\subsection{本章小结}
Ad hoc teamwork 把多 Agent 问题从 ``是否能协同'' 推进到 ``是否能和未知协同体稳定协同''，这是落地系统必须面对的版本。

